\chapter{Premiers pas en PHP}

\begin{synthese}[Au programme]
Insérer du PHP dans une page HTML, comprendre où s'exécute un script, déclarer
variables et constantes, utiliser les opérateurs et afficher des données avec
\code{echo} et \code{print}. C'est le socle de toute application web dynamique.
\end{synthese}

\section{Qu'est-ce que PHP ?}

\begin{definition}
\textbf{PHP} (\emph{PHP: Hypertext Preprocessor}) est un langage de script
\textbf{exécuté côté serveur}. Le serveur interprète le code PHP, puis envoie au
navigateur du client une page \textbf{HTML} ordinaire : le visiteur ne voit jamais
le code PHP.
\end{definition}

\begin{remarque}
Contrairement à JavaScript qui s'exécute dans le navigateur (côté \emph{client}),
PHP s'exécute sur le \emph{serveur}. C'est pour cela qu'on peut, avec PHP, se
connecter à une base de données ou lire des fichiers du serveur en toute sécurité.
\end{remarque}

\section{Insérer du PHP dans une page}

Un script PHP se place entre les balises \code{<?php} et \code{?>}. Le fichier doit
porter l'extension \code{.php}.

\begin{syntaxe}
\code{<?php}\quad \textit{instructions PHP\dots}\quad \code{?>}
\end{syntaxe}

\begin{codephp}
<!DOCTYPE html>
<html lang="fr">
<head><meta charset="utf-8"><title>Ma première page</title></head>
<body>
  <h1>Bonjour</h1>
  <?php
    // Ceci est un commentaire sur une ligne
    echo "<p>Bienvenue sur OnBuch, il est " . date("H:i") . " !</p>";
  ?>
</body>
</html>
\end{codephp}

\fig{Le code PHP est noyé dans le HTML ; seul son résultat est envoyé au client.}

\begin{attention}
N'oublie jamais le point-virgule \code{;} à la fin de chaque instruction PHP, ni la
balise fermante \code{?>}. Un oubli provoque une erreur de syntaxe (\emph{Parse error}).
\end{attention}

\section{Variables et constantes}

\begin{definition}
Une \textbf{variable} est un espace mémoire nommé qui stocke une valeur modifiable.
En PHP, le nom d'une variable commence \textbf{toujours} par le symbole \code{\$}.
\end{definition}

\begin{codephp}
<?php
  $nom = "Awa";          // chaîne de caractères
  $age = 17;              // entier
  $moyenne = 14.5;        // nombre décimal
  $estAdmis = true;       // booléen
  define("ETABLISSEMENT", "Lycée de Maroua");  // constante
  echo "Élève : $nom ($age ans) — " . ETABLISSEMENT;
?>
\end{codephp}

\begin{propriete}[Règles de nommage]
Un nom de variable commence par \code{\$} suivi d'une lettre ou d'un \code{\_},
puis de lettres, chiffres ou \code{\_}. PHP distingue les majuscules des minuscules :
\code{\$nom} et \code{\$Nom} sont deux variables différentes.
\end{propriete}

\section{Afficher : \texttt{echo} et \texttt{print}}

\begin{itemize}
  \item \code{echo} affiche une ou plusieurs valeurs (séparées par des virgules) ;
  \item \code{print} affiche une seule valeur et renvoie \code{1} ;
  \item l'opérateur de concaténation \code{.} colle deux chaînes.
\end{itemize}

\begin{exemple}
\code{echo "Total : ", 3 + 4;} affiche \code{Total : 7}, tandis que
\code{echo "3 + 4";} affiche littéralement \code{3 + 4}.
\end{exemple}

\section{Les opérateurs}

\begin{center}\small
\begin{tabular}{@{}lll@{}}
\toprule
\textbf{Catégorie} & \textbf{Opérateurs} & \textbf{Exemple}\\
\midrule
Arithmétiques & \code{+ - * / \% **} & \code{\$x \% 2} (reste)\\
Comparaison & \code{== != < > <= >= ===} & \code{\$a === \$b}\\
Logiques & \code{\&\& || !} & \code{\$x>0 \&\& \$x<10}\\
Affectation & \code{= += -= *= .=} & \code{\$s .= "!"}\\
\bottomrule
\end{tabular}
\end{center}

\begin{methode}
\textbf{Tester un script PHP en local.} (1)~Lancer le serveur (WampServer/XAMPP) ;
(2)~placer le fichier \code{.php} dans le dossier \code{www} (ou \code{htdocs}) ;
(3)~ouvrir \code{http://localhost/monfichier.php} dans le navigateur.
\end{methode}

\section{Exercices}

\rubrique{Application}

\exo{Première page dynamique}\diff{1}
Écris un script PHP qui affiche dans une page HTML le message
\emph{« Nous sommes le \dots »} suivi de la date du jour.

\exo{Calculs et affichage}\diff{1}
On donne \code{\$a = 7} et \code{\$b = 3}. Affiche, chacune sur une ligne, la somme,
la différence, le produit, le quotient et le reste de la division de \code{\$a} par \code{\$b}.

\exo{Bulletin}\diff{2}
Trois notes sont stockées dans \code{\$n1}, \code{\$n2}, \code{\$n3}. Calcule et affiche
la moyenne avec deux décimales (fonction \code{round}).

\section{Corrigés}

\corrige{de l'exercice 1}
\begin{codephp}
<?php
  echo "<p>Nous sommes le " . date("d/m/Y") . "</p>";
?>
\end{codephp}

\corrige{de l'exercice 2}
\begin{codephp}
<?php
  $a = 7; $b = 3;
  echo "Somme : ", $a + $b, "<br>";
  echo "Différence : ", $a - $b, "<br>";
  echo "Produit : ", $a * $b, "<br>";
  echo "Quotient : ", $a / $b, "<br>";
  echo "Reste : ", $a % $b;
?>
\end{codephp}

\corrige{de l'exercice 3}
\begin{codephp}
<?php
  $n1 = 12; $n2 = 15; $n3 = 9;
  $moyenne = ($n1 + $n2 + $n3) / 3;
  echo "Moyenne : " . round($moyenne, 2);
?>
\end{codephp}

\begin{aretenir}
PHP s'exécute \textbf{côté serveur} entre \code{<?php} et \code{?>}. Les variables
commencent par \code{\$}, chaque instruction finit par \code{;}, et \code{echo}
sert à renvoyer du HTML au navigateur.
\end{aretenir}
